\chapter{Internship Context Presentation}
\label{ch:context}

This chapter establishes the context of the internship, introduces the
host research institution, and sets out the objectives, work packages and
expected outcomes of the project.

\section{Presentation of the Host Institution}

\subsection{Overview}

The internship was carried out at the \textbf{Maersk Mc-Kinney Moller
Institute}, a research institute of the \textbf{University of Southern
Denmark (SDU)}. The Institute is an academic research environment
organised around robotics, artificial intelligence, embedded systems and
applied data science, and its work is structured around research
questions and publications rather than around product delivery for
clients.

This distinction shaped the entire internship. There was no client, no
specification document and no delivery deadline for a feature set. The
deliverable was a defensible answer to a research question, together
with the experimental apparatus needed to reproduce it --- which, as
Chapter~\ref{ch:results} discusses, includes the apparatus needed to
\emph{disprove} an initially promising result.

\subsection{Internship Framework}

This work fulfils the research internship requirement of the first-year
engineering cycle at ISSAT Sousse. The internship was conducted
\textbf{remotely}, under the academic supervision of \textbf{Professor
Sadok Ben Yahia}, over the period 15 June 2026 to 15 August 2026.

The remote format had a concrete methodological consequence. All
computation ran as batch jobs on a high-performance computing cluster
under the SLURM scheduler, with no interactive session and no ability to
attach a debugger to a running job. Every experiment therefore had to be
specified completely in advance, instrumented well enough to be
diagnosed from its log file alone, and cheap enough to fail fast. This
constraint is visible throughout the engineering choices described in
Chapter~\ref{ch:tasks}.

\section{Context and Objectives of the Internship}

\subsection{Purpose}

Explainable Artificial Intelligence (XAI) for medical imaging is
dominated by attribution methods, which localise evidence but do not
describe the change that would alter a decision. Counterfactual image
generation supplies that missing description, but it introduces a
verification problem of its own: an image that flips a classifier's
output has not necessarily acquired --- or lost --- any disease. It may
simply have acquired an adversarial perturbation tailored to that one
model.

The purpose of this project was therefore twofold:

\begin{itemize}[leftmargin=*]
  \item To build a counterfactual generation system for binary pneumonia
        classification on chest radiographs, in which the diffusion
        sampling schedule is \emph{learned} by a reinforcement-learning
        agent rather than fixed by hand.
  \item To evaluate that system honestly --- including with controls
        capable of exposing the adversarial failure mode --- rather than
        reporting only the metric that the system was optimised against.
\end{itemize}

\subsection{Main Objective}

The central objective was to engineer and evaluate a complete research
pipeline that, given a chest radiograph $x$ and a target class $c$,
produces a counterfactual $x_c$ satisfying three simultaneous criteria:

\begin{enumerate}[leftmargin=*]
  \item \textbf{Validity}: a frozen classifier $f$ scores $x_c$ at the
        target class with a large margin.
  \item \textbf{Realism}: $x_c$ is distributionally indistinguishable
        from genuine radiographs.
  \item \textbf{Identity preservation}: $x_c$ differs from $x$ only in a
        minimal, class-relevant region, leaving patient anatomy and
        foreign objects intact.
\end{enumerate}

The methodological commitment underlying the whole system is that the
diffusion model's \emph{weights} are never optimised by reinforcement
learning. The PPO agent controls only the sampling procedure --- noise
strength, guidance scales, number of refinement passes, and the stopping
decision --- around a frozen, LoRA-adapted generator. This keeps the
action space small and interpretable, and makes the learned policy
reusable across arbitrary query images at inference time.

\subsection{Assigned Work Packages}

The project was decomposed into six work packages:

\begin{itemize}[leftmargin=*]
  \item \textbf{WP1 --- Data ingestion.} Programmatic acquisition of the
        RSNA Pneumonia Detection Challenge dataset, DICOM decoding,
        normalisation, and derivation of a binary label from the
        bounding-box annotations.
  \item \textbf{WP2 --- Pretrained model integration.} Loading and
        freezing a pneumonia classifier, a lung/heart segmentation
        network, and an object-detector feature extractor, all serving as
        fixed sources of guidance and reward signal.
  \item \textbf{WP3 --- Generator adaptation.} Parameter-efficient
        LoRA fine-tuning of a Stable Diffusion~1.5 UNet from the
        natural-image domain to chest radiography.
  \item \textbf{WP4 --- Reinforcement-learning control.} Design of a
        Gym-style environment, a shaped reward function, and a PPO
        actor-critic agent; training of two directional policies, one per
        class-flip direction.
  \item \textbf{WP5 --- Evaluation.} Implementation of the
        Singla~\emph{et al.}~\cite{singla2022} metric suite (CV, FID,
        FOP) plus clinical proxy metrics, and of pixel-level difference
        visualisation.
  \item \textbf{WP6 --- Validation and diagnostics.} Independent
        held-out re-scoring to distinguish genuine disease signal from
        adversarial perturbation, and Grad-CAM analysis of the guidance
        classifier for shortcut learning.
\end{itemize}

WP6 was not part of the initial plan. It was added mid-internship, in
direct response to the observation that the primary validity metric was
computed with the same classifier that had produced the guidance
gradient, and was therefore circular. Its results are the central finding
of this report.

A further piece of work, undertaken after WP1--WP6 and reported separately
in Chapter~\ref{ch:breast}, was not assigned and is not counted among these
six packages: a self-initiated extension of the pipeline to a second
imaging domain, breast-cancer screening mammography, to test whether the
architecture and the WP6 evaluation discipline generalise beyond chest
radiography.

\section{Expected Results}

The anticipated deliverables at the outset of the internship were:

\begin{itemize}[leftmargin=*]
  \item A reproducible, end-to-end research pipeline covering ingestion,
        generator adaptation, policy training and evaluation, executable
        as batch jobs on an HPC cluster.
  \item Two trained directional policies together with their
        checkpoints and configuration provenance.
  \item A quantitative evaluation under the reference metric protocol,
        with the numbers traceable to serialised artefacts rather than
        to transcription.
  \item A critical assessment of what the system does and does not
        demonstrate, including its failure modes.
\end{itemize}

\section{Conclusion}

This chapter has set out the academic research context of the
internship, introduced the Maersk Mc-Kinney Moller Institute as host
institution, and defined the objectives and work packages of the
project. The following chapter details the research and development
environment --- the cluster, the frameworks, and the pretrained model
zoo --- within which this work was carried out.
